\section{High Fidelity Enthalpy\label{sec:highFidelity}}
\subsection{\acrfull{StS}\label{sec:sts}}
\gls{StS} simulations were run for a range of conditions relevant to hypersonic
flight conditions for pure oxygen. This included a freestream typical of
$60\Units{km}$ altitude at different Mach numbers $(5 - 8)$. The vibrational
population distribution evolution post-normal shock waves at $60\Units{km}$ is
shown in \Cref{fig:stsDistribution} for Mach $5,6$ and $8$ as a function of
time behind a normal shock using \gls{QCT} based \gls{StS} simulations.
Significant differences can be seen at higher vibrational states but these have
a significantly lower population concentration compared to the lowest energy
states.
\begin{figure}[H]
    \centering
    \subfigure[Mach $5$]{\label{fig:stsDistributionM5}\includegraphics[width=0.45\textwidth]{chapter4/stsStudy/distributionMach5.pdf}}
    \subfigure[Mach
    $6$]{\label{fig:stsDistributionM6}\includegraphics[width=0.45\textwidth]{chapter4/stsStudy/distributionMach6.pdf}}
    \subfigure[Mach
    $8$]{\label{fig:stsDistributionM8}\includegraphics[width=0.45\textwidth]{chapter4/stsStudy/distributionMach8.pdf}}
    \caption{Vibrational population distribution evolution post-normal shock
    waves at $60\Units{km}$\label{fig:stsDistribution}}
\end{figure}
\Cref{fig:stsDistribution} illustrates the differences between the
non-Boltzmann vibrational energy distribution computed using the \gls{StS}
model and the classical Boltzmann distribution. The Boltzmann distribution
assumes thermal equilibrium among vibrational states, whereas the \gls{StS}
model accounts for nonequilibrium effects and does not assume thermal
equilibrium.\\

The discrepancy between the \gls{StS} and Boltzmann distributions is most
pronounced at very short times after the initial condition. As time progresses,
the distributions tend to converge for lower vibrational quantum numbers,
indicating relaxation toward equilibrium. Notably, the highest Mach number case
shown in \Cref{fig:stsDistributionM8} achieves a steady-state distribution more
rapidly (around $1,\mu s$) compared to the lower Mach number cases in
\Cref{fig:stsDistributionM5} and \Cref{fig:stsDistributionM6}. This faster
relaxation is consistent with the increased energy and collision rates at
higher flow velocities. \\

Using the results from the \gls{StS} simulations and the \textit{Buldakov
approach} \cite{Buldakov2000}, the vibrational-state specific polarizability
was calculated. These results are presented in \Cref{fig:stsTemperatures}.
Three important trends can be observed. First, the polarizability exhibits a
strong correlation with the vibrational temperature and increases with time
downstream of the normal shock. This behavior is significant since many optical
sensors used for flow diagnostics are positioned downstream of the stagnation
point, where vibrational excitation is prominent. Second, as time progresses,
the difference in polarizability between the two distributions (Boltzmann and
non-Boltzmann) increases for each Mach number case before reaching a plateau,
indicating a steady-state behavior in the nonequilibrium population dynamics.
Third, the Mach 8 case demonstrates the largest discrepancy between the two
distribution types, highlighting the enhanced nonequilibrium effects at higher
Mach numbers.
\begin{figure}[H]
    \centering
    \subfigure[Mach $5$]{\label{fig:stsTemperatureM5}\includegraphics[width=0.45\textwidth]{chapter4/stsStudy/temperatureMach5.pdf}}
    \subfigure[Mach
    $6$]{\label{fig:stsTemperatureM6}\includegraphics[width=0.45\textwidth]{chapter4/stsStudy/temperatureMach6.pdf}}
    \subfigure[Mach
    $8$]{\label{fig:stsTemperatureM8}\includegraphics[width=0.45\textwidth]{chapter4/stsStudy/temperatureMach8.pdf}}
    \caption{Temperature profiles and polarizability profiles post normal shock
    waves at $60\Units{km}$\label{fig:stsTemperatures}}
\end{figure}
Although the polarizability shows a strong dependence on vibrational
temperature, its dependence on the specific vibrational distribution (Boltzmann
versus non-Boltzmann) is comparatively weaker. This is because the molecular
population remains predominantly concentrated in the lower vibrational energy
levels for both distributions. While a more pronounced difference between the
two distributions appears at higher vibrational energy levels, the population
in these higher states is orders of magnitude smaller. Consequently, their
contribution to the overall polarizability is limited. This suggests that for
practical purposes, vibrational temperature is the primary factor influencing
polarizability, with deviations from Boltzmann equilibrium having a secondary
effect, particularly in flows where lower vibrational states dominate the
molecular population.
